\section{Riemannian Geometry}
\label{sec:riemannian_geometry}

Now that we've covered the basics of differential geometry (cf Section~\ref{sec:differential_geometry}), we now give a brief introduction to Riemannian Geometry.

While differential geometry gives us the structure of the manifold, \emph{Riemannian geometry} introduces the tools necessary for measurement: lengths, angles, and volumes. Riemannian geometry coincides with our intuitive understanding that geometry is the study of lengths, angles, curvature, area, volume, and so on. It provides the mathematical framework for understanding the geometry of curved spaces, which is essential for many applications in physics, engineering, and other fields.

\subsection{The Metric Tensor} Let $M$ be a manifold. We can define something called a metric that allows us to talk about distances and angles on the manifold. We can define a metric as follows.

\begin{definition}
  A \emph{metric tensor} on $M$ is a smooth assignment
  \[%
    g_p : T_pM \times T_pM \to \R
  ,\]%
  for each point $p \in M$ such that:
  \begin{enumerate}
    \item $g_p$ is bilinear,
    \item $g_p$ is symmetric,
    \item $g_p$ is nondegenerate.
  \end{enumerate}
\end{definition}

Smoothness means that in any coordinate chart $\{x^\mu\}$, the components
\[%
  g_{\mu\nu}(x) := g\!\left(\pd{\mu}, \pd{\nu}\right)
,\]%
are smooth functions. From this, we get a special type of manifold, namely, the pseudo-Riemannian manifold.

\begin{definition}
  The pair $(M, g)$ is called a \emph{pseudo-Riemannian manifold}.

  If in addition $g_p$ is positive definite for every $p \in M$, meaning
  \[%
    g_p(v,v) > 0 \quad \text{for all}~v \neq 0
  ,\]%
  then $(M, g)$ is called a \emph{Riemannian manifold}.
\end{definition}

In general, a pseudo-Riemannian metric need not be positive definite. At each point $p$, the bilinear form $g_p$ has a well-defined \emph{signature} $(p, q)$, where $p$ and $q$ denote the number of positive and negative eigenvalues, respectively.

The metric allows us to:

\begin{itemize}
  \item Measure lengths of tangent vectors:
    \[%
      \|v\|^2 = g(v,v)
    ,\]%
  \item Measure angles between vectors,
  \item Raise and lower tensor indices using the inverse metric $g^{\mu\nu}$.
\end{itemize}

Because $g$ is nondegenerate, there exists a unique inverse tensor $g^{\mu\nu}$ satisfying
\[%
  g^{\mu\alpha} g_{\alpha\nu} = \delta^\mu_{\nu}
.\]%

This allows us to lower and raise indices:
\[%
  v_\mu = g_{\mu\nu} v^\nu \aand v^\mu = g^{\mu\nu} v_\nu
.\]%

Metrics that are not positive definite naturally classify tangent vectors into different types depending on the sign of $g(v, v)$. A particularly important case is when the signature is $(1, n-1)$, which will be discussed in detail later when we introduce Lorentzian manifolds.

\subsection{Levi-Civita Connection}

Given a pseudo-Riemannian manifold $(M, g)$, there is a unique connection that is compatible with the metric and has zero torsion. This connection is called the \emph{Levi-Civita connection}.

\begin{theorem}[Fundamental Theorem of Riemannian Geometry]
  Let $(M, g)$ be a pseudo-Riemannian manifold. There exists a unique connection $\nabla$ (the \emph{Levi-Civita connection}) such that:
  \begin{enumerate}
    \item $\nabla$ is torsion-free: $\nabla_X Y - \nabla_Y X = [X,Y]$.
    \item $\nabla$ is metric-compatible: $X(g(Y,Z)) = g(\nabla_X Y, Z) + g(Y, \nabla_X Z)$.
  \end{enumerate}
\end{theorem}

From the properties of the Levi-Civita connection, we can derive an explicit formula for $\nabla$ in terms of the metric, known as the \emph{Koszul formula}:
\[%
  2 g(\nabla_X Y, Z) = X(g(Y,Z)) + Y(g(Z,X)) - Z(g(X,Y)) + g([X,Y], Z) - g([Y,Z], X) - g([Z,X], Y)
.\]%
In a local coordinate system $\{x^i\}$, we can express the covariant derivative of one basis vector $\pd{j}$ in the direction of another $\pd{i}$. Since the result $\nabla_{\pd{i}} \pd{j}$ is itself a vector field, it can be written as a linear combination of the basis vectors $\{\pd{k}\}$.

\begin{definition}
  The \emph{Christoffel symbols of the second kind} are defined by the relation:
  \[%
    \nabla_{\pd{i}} \pd{j} = \Gamma^k_{ij} \pd{k}
  .\]%

  Applying the Koszul formula to the coordinate basis vectors (where $[\pd{i}, \pd{j}] = 0$) yields the explicit formula:
  \[%
    \Gamma^k_{ij} = \frac{1}{2} g^{kl} \left( \pd{i} g_{jl} + \pd{j} g_{il} - \pd{l} g_{ij} \right)
  ,\]%
  where $g^{kl}$ is the inverse of the metric tensor $g_{kl}$. This formula allows us to compute the Christoffel symbols directly from the metric components and their derivatives.
\end{definition}

Since the Levi-Civita connection is torsion-free, the Christoffel symbols are symmetric in the lower indices:
\[%
  \Gamma^k_{ij} = \Gamma^k_{ji}
.\]%
This reduces the number of components you have to calculate (e.g., in 3D, there are only 18 independent components instead of 27).

Now, let's use the Christoffel symbols to compute the covariant derivative of a vector field $V = V^k \pd{k}$ in the direction of another vector field $X = X^i \pd{i}$. Let $V = V^k\pd{k}$ be a smooth vector field and $X = X^i\pd{i}$ be the direction of differentiation. By the Leibniz rule (property of the connection), the covariant derivative $\nabla_X V$ is:
\[%
  \nabla_XV = \nabla_{X^i\pd{i}}(V^j\pd{j}) = X^i\nabla_{\pd{i}}(V^j\pd{j})
.\]%
Applying the product rule:
\[%
  \nabla_XV = X^i\left(\left(\pd{i}V^j\right)\pd{j} + V^j\nabla_{\pd{i}}\pd{j}\right)
.\]%
Recalling the definition of Christoffel symbols, $\nabla_{\pd{i}}\pd{j} = \Gamma_{ij}^k \pd{k}$, we substitute this into the expression:
\[%
  \nabla_XV = X^i\left(\left(\pd{i}V^k\right)\pd{k} + V^j\Gamma_{ij}^k \pd{k}\right)
.\]%
Factoring out the basis $\pd{k}$, we arrive at the component formula:
\[%
  (\nabla_X V)^k = X^i \left(\pdv{V^k}{x^i} + \Gamma^k_{ij} V^j\right)
.\]%

\subsection{Parallel Transport and Geodesics} In Euclidean space, a ``straight'' path is easily defined as one with zero acceleration. However, on a Riemannian manifold $(M, g)$, the notion of a straight line is complicated by the fact that the tangent spaces $T_p M$ and $T_q M$ at different points are distinct. To determine if a path is straight, we first need a way to move vectors from one point to another without ``turning'' them. This leads us to the concept of parallel transport.

Let $(M, g)$ be a manifold equipped with the Levi-Civita connection $\nabla$. Given a smooth curve $\gamma: I \to M$ with tangent vector $\dot{\gamma}(t) = \frac{d\gamma}{dt}$, we seek to measure how a vector field $X$ changes as we move along $\gamma$.

\begin{definition}
  Let $X$ be a vector field defined along $\gamma$, such that $X(t) \in T_{\gamma(t)} M$. The \emph{covariant derivative of $X$ along $\gamma$} is the projection of the standard derivative back onto the tangent space:
  \[%
    \codv{X}{t} := \nabla_{\dot{\gamma}(t)} X
  .\]%
\end{definition}

In local coordinates $\{x^\mu\}$, the covariant derivative accounts for the ``twisting'' of the coordinate basis through the Christoffel symbols $\Gamma^\mu_{\alpha\beta}$:
\[%
  \left(\codv{X}{t}\right)^\mu = \odv{X^\mu}{t} + \Gamma^\mu_{\alpha\beta} \dot{\gamma}^\alpha X^\beta
.\]%
This formula captures the motivation for the connection: it corrects the ordinary derivative of the component $X^\mu$ by adding a term that represents how the basis vectors themselves are changing relative to the metric.

A vector is ``parallel'' if it does not change its orientation relative to the manifold's geometry as it moves.

\begin{definition}
  A vector field $X$ along $\gamma$ is said to be \emph{parallel transported} if its covariant derivative along the curve vanishes:
  \[%
    \codv{X}{t} = 0 \quad \iff \quad \odv{X^\mu}{t} + \Gamma^\mu_{\alpha\beta} \dot{\gamma}^\alpha X^\beta = 0
  .\]%
\end{definition}

Because the Levi-Civita connection is metric-compatible, parallel transport is an \emph{isometry} between tangent spaces. It preserves both the lengths of vectors and the angles between them:
\[%
  \odv{}{t} g(X, Y) = g\left(\codv{X}{t}, Y\right) + g\left(X, \codv{Y}{t}\right) = 0
.\]%
In medical imaging, this ensures that if an ultrasound pulse is modeled as a vector, its ``energy'' (magnitude) remains constant as it moves through the tissue, even if its direction is deflected by the curvature of the sound-speed metric.

With parallel transport defined, we can now define the ``straightest'' possible curves on a manifold. A curve is a geodesic if it does not ``accelerate'' in any direction other than the one dictated by the curvature of the manifold itself.

\begin{definition}
  A curve $\gamma$ is a \emph{geodesic} if its tangent vector is parallel transported along itself:
  \[%
    \codv{\dot{\gamma}}{t} = 0 \quad \text{or equivalently,} \quad \nabla_{\dot{\gamma}} \dot{\gamma} = 0
  .\]%
\end{definition}

Expanding this in local coordinates yields the \emph{geodesic equation}, a second-order non-linear system of ODEs:
\[%
  \odv[2]{x^\mu}{t} + \Gamma^\mu_{\alpha\beta} \odv{x^\alpha}{t} \odv{x^\beta}{t} = 0
.\]%

While the definition above is geometric (zero acceleration), geodesics also arise from a physical principle: the minimization of effort. Geodesics are critical points of the \emph{energy functional}:
\[%
  E[\gamma] = \frac{1}{2} \int_I g_{\mu\nu} \dot{x}^\mu \dot{x}^\nu \dt
.\]%
Applying the Euler-Lagrange equations to $E[\gamma]$ yields the exact same coordinate geodesic equation. In the context of our boundary rigidity problem, this variational perspective is vital; it tells us that ultrasound rays (geodesics) always follow the paths that minimize ``acoustic effort'' (time-of-flight), which is the very data we measure at the boundary.

\subsection{Curvature} Next, we introduce the concept of curvature on a manifold.

\begin{definition}
  The \emph{Riemann curvature tensor} is the $(1,3)$-tensor $R: \Gamma(TM) \times \Gamma(TM) \times \Gamma(TM) \to \Gamma(TM)$ defined by:
  \begin{equation}
    R(X,Y)Z = \nabla_X \nabla_Y Z - \nabla_Y \nabla_X Z - \nabla_{[X,Y]} Z
  .\end{equation}
\end{definition}

This tensor quantifies the failure of the second covariant derivatives to commute. Geometrically, it measures the change in a vector $Z$ after it is parallel-transported around an infinitesimal loop spanned by $X$ and $Y$.

In a local coordinate system $\{x^\mu\}$, the components of the Riemann tensor are computed from the Christoffel symbols and their partial derivatives:
\[%
  R^\rho_{\sigma\mu\nu} = \pd{\mu} \Gamma^\rho_{\nu\sigma} - \pd{\nu} \Gamma^\rho_{\mu\sigma} + \Gamma^\rho_{\mu\lambda} \Gamma^\lambda_{\nu\sigma} - \Gamma^\rho_{\nu\lambda} \Gamma^\lambda_{\mu\sigma}
.\]%

The Riemann tensor for a Levi-Civita connection satisfies several symmetries, most notably $R(X,Y)Z = -R(Y,X)Z$, which in coordinates reads $R^\rho_{\sigma\mu\nu} = -R^\rho_{\sigma\nu\mu}$. By contracting indices, we obtain the primary tools for geometric analysis:
\begin{itemize}
  \item \textbf{Ricci Tensor:} $R_{\sigma\nu} = R^\rho_{\sigma\rho\nu}$. This $(0,2)$-tensor represents the average curvature in a given direction.
  \item \textbf{Scalar Curvature:} $R = g^{\sigma\nu} R_{\sigma\nu}$. This is a single smooth function providing a global measure of how the manifold's volume differs from Euclidean space.
\end{itemize}

\begin{example}
  For a 2-sphere $\S^2$ of radius $a$ with the standard metric $\ds^2 = a^2 \dd{\theta}^2 + a^2 \sin^2 \theta \dd{\phi}^2$, a direct computation of the Riemann components reveals that the curvature is constant. Specifically, the scalar curvature is:
  \[%
    R = \frac{2}{a^2}
  .\]%
  In medical imaging, regions with high $R$ might correspond to areas of high structural complexity or sharp anatomical transitions.
\end{example}

If a metric $g_{\mu\nu}$ is given, we may further contract to obtain the scalar curvature.

\begin{definition}
  The \emph{scalar curvature} is defined by
  \[%
    R = g^{\sigma\nu} R_{\sigma\nu}
  \]%
\end{definition}

While the Riemann tensor is a high-rank object, we can extract the curvature of a specific 2-dimensional plane $P \subset T_pM$ within the tangent space. This is often the most intuitive way to visualize curvature in higher dimensions.
\begin{definition}
  Let $\{v, w\}$ be a basis for a 2-plane $P$ in $T_pM$. The \emph{sectional curvature} of $P$ is defined by:
  \[%
    K(v,w) = \frac{g(R(v,w)w, v)}{g(v,v)g(w,w) - g(v,w)^2}
  .\]%
\end{definition}
For a Riemannian manifold, if $K$ is constant for all planes at all points, the manifold is called a \emph{space of constant curvature}. In imaging, sectional curvature helps characterize the ``local shape'' of an organ's surface.

In dimensions $n \geq 4$, the Riemann tensor can be decomposed into a part determined by the Ricci tensor and a ``remainder'' known as the Weyl tensor.
\begin{definition}
  The \emph{Weyl curvature tensor} $C$ is the traceless part of the Riemann tensor. It represents the component of curvature that distorts the shape of a volume element without changing its total volume.
\end{definition}
This is particularly relevant for \emph{conformal geometry}, where we study deformations that preserve angles but not necessarily distances--a common scenario when modeling tissue stretching.

Finally, for variational problems and physical modeling, we define a specific combination of the Ricci and Scalar curvatures.
\begin{definition}
  The \emph{Einstein tensor} $G$ is the symmetric $(0,2)$-tensor defined by:
  \[%
    G_{\mu\nu} = R_{\mu\nu} - \frac{1}{2} R g_{\mu\nu}
  .\]%
\end{definition}
The primary mathematical utility of $G$ is that its divergence vanishes ($\nabla^\mu G_{\mu\nu} = 0$) due to the \emph{contracted Bianchi identity}. This property is essential when setting up ``conservation laws'' for biological flows or deformations.

\subsection{Curvature Contractions}

The Riemann tensor possesses 256 components in four dimensions, but its inherent symmetries reduce this to only 20 independent components. These symmetries are essential for performing contractions correctly.

\begin{itemize}
  \item \textbf{Algebraic Symmetries:}
    \begin{align}
      R_{\rho\sigma\mu\nu} &= -R_{\sigma\rho\mu\nu} \quad \text{(Antisymmetry in first two indices)} \\
      R_{\rho\sigma\mu\nu} &= -R_{\rho\sigma\nu\mu} \quad \text{(Antisymmetry in last two indices)} \\
      R_{\rho\sigma\mu\nu} &= R_{\mu\nu\rho\sigma} \quad \text{(Interchange symmetry)}
    \end{align}

  \item \textbf{The First Bianchi Identity:} This relates the cyclic permutation of the last three indices:
    \[ R^\rho_{\sigma\mu\nu} + R^\rho_{\mu\nu\sigma} + R^\rho_{\nu\sigma\mu} = 0 \]

  \item \textbf{The Second (Differential) Bianchi Identity:} This describes how the curvature tensor varies across the manifold:
    \[ \nabla_\lambda R^\rho_{\sigma\mu\nu} + \nabla_\mu R^\rho_{\sigma\nu\lambda} + \nabla_\nu R^\rho_{\sigma\lambda\mu} = 0 \]
\end{itemize}

By contracting the second Bianchi identity twice, we arrive at a result that is fundamental to variational principles in imaging (such as Total Variation denoising on manifolds):
\begin{equation}
  \nabla^\mu R_{\mu\nu} = \frac{1}{2} \nabla_\nu R
.\end{equation}
This identity implies that the \emph{Einstein tensor} $G_{\mu\nu} = R_{\mu\nu} - \frac{1}{2}Rg_{\mu\nu}$ is divergence-free ($\nabla^\mu G_{\mu\nu} = 0$). In physical modeling, this represents a local conservation law for the geometric ``energy'' of the manifold.

The Riemann tensor can be decomposed into three irreducible components that represent different geometric effects. This decomposition is particularly useful in higher-dimensional analysis (where $n \geq 3$).
\begin{enumerate}
  \item \textbf{The Scalar Part:} Represents the portion of curvature associated with the scalar curvature $R$. This part determines how the volume of a small ball in the manifold compares to a ball of the same radius in Euclidean space.

  \item \textbf{The Traceless Ricci Part:} Encapsulated by the tensor $R_{\mu\nu} - \frac{1}{n} R g_{\mu\nu}$. This describes the ``tidal'' forces or the anisotropic stretching of the volume that preserves the total volume but changes its shape.

  \item \textbf{The Weyl Tensor ($C_{\rho\sigma\mu\nu}$):} Also known as the \emph{conformal curvature tensor}. In dimensions $n \geq 4$, it is the part of the Riemann tensor that remains after the Ricci and Scalar parts are removed.
\end{enumerate}

To connect these contractions to the ``Inverse Problem,'' consider a small ``cloud'' of test particles (or signal points in an MRI) initially at rest. The Ricci tensor $R_{\mu\nu}$ describes the initial second derivative of the volume $V$ of this cloud as it begins to move along geodesics:
\begin{equation}
  \left. \frac{\ddot{V}}{V} \right\rvert_{t=0} = -R_{\mu\nu} \dot{\gamma}^\mu \dot{\gamma}^\nu
.\end{equation}
This relation, known as \emph{Raychaudhuri's equation} in its simplest form, shows that positive Ricci curvature tends to converge geodesics, causing the volume to shrink. In imaging, this can be used to model the focus or dispersion of signals within heterogeneous media.

\subsection{The Laplace-Beltrami Operator}

In Euclidean space, the Laplacian $\Delta = \sum \pd{i}^2$ is the most fundamental second-order differential operator. On a Riemannian manifold $(M, g)$, this is generalized to the \emph{Laplace-Beltrami operator} $\Delta_g$, which accounts for the varying metric and curvature of the space.
\begin{definition}
  The \emph{Laplace-Beltrami operator} acting on a smooth function $f \in C^\infty(M)$ is defined as the divergence of the gradient:
  \begin{equation}
    \Delta_g f = \div(\grad f) = \nabla^\mu \nabla_\mu f
  .\end{equation}
\end{definition}

Using the metric $g$ and its determinant $|g| = \det(g_{\mu\nu})$, the Laplace-Beltrami operator can be written explicitly without requiring the Christoffel symbols:
\begin{equation}
  \Delta_g f = \frac{1}{\sqrt{|g|}} \pd{\mu} \left( \sqrt{|g|} g^{\mu\nu} \pd{\nu} f \right)
.\end{equation}
This formula is particularly powerful for medical imaging simulations because it only requires the metric components $g_{\mu\nu}$ and their first derivatives.

The study of the equation $\Delta_g \phi = -\lambda \phi$ is known as the \emph{spectral geometry} of the manifold. For a compact manifold, the eigenvalues $0 = \lambda_0 < \lambda_1 \leq \lambda_2 \dots$ form a discrete spectrum that encodes global geometric information, such as the volume and total curvature (by the Heat Kernel expansion).

In the context of medical imgaging, the Laplace-Beltrami operator appears in several key areas:
\begin{itemize}
  \item \textbf{Diffusion Tensor Imaging (DTI):} The diffusion of water molecules in the brain follows the heat equation $\pd{t} u = \Delta_g u$, where $g$ is the diffusion tensor.
  \item \textbf{Surface Parameterization:} The Laplace-Beltrami operator is used to calculate ``Harmonic maps,'' which allow researchers to flatten complex anatomical surfaces (like the cerebral cortex) into 2D planes while minimizing distortion.
  \item \textbf{Geometric Regularization:} In inverse problems, the ``energy'' term $\int_M |\nabla_g f|^2 dV$ is often used to ensure that reconstructed images are smooth relative to the underlying anatomy.
\end{itemize}
